\section{A necessary gate for an honest CY3 realization}
\label{sec:realization}

The cubic in \cref{thm:yukawa} is an invariant of the pointed polarized
variation, up to projective tangent change and a common scalar.  It
therefore supplies a one-sided realization obstruction.

\begin{theorem}[Projective Yukawa necessity]
\label{thm:realization-gate}
Let \(g:Y\to T\) be a smooth projective family of Calabi--Yau threefolds
with \(h^{2,1}=4\).  Suppose there are an isomorphism of pointed complex
base germs
\[
 \phi:(T,t)\xrightarrow{\sim}(\Bcore,0)
\]
and an isomorphism of polarized rational variations of Hodge structure
\[
 \Phi:R^3g_*\Q\xrightarrow{\sim}\phi^*\Vcore.       \tag{8.1}
\label{eq:8.1}
\]
If \(A=d\phi_t\), then the two Yukawa cubics satisfy
\[
 Y_Y(v)=\lambda\,Y_H(Av)
 \quad\text{for some }\lambda\in\CC^\times.         \tag{8.2}
\label{eq:8.2}
\]
Thus failure of projective \(\operatorname{GL}_4(\CC)\)-equivalence is a
pointed local polarized-VHS obstruction.
\end{theorem}

\begin{proof}
The isomorphism \(\Phi\) is horizontal, filtration-preserving, and
polarization-preserving.  It identifies the one-dimensional \(F^3\) lines.
Choose local generators \(\Omega_Y\) and \(\Omega_H\); at the base point
their images differ by a nonzero scalar.  Horizontality intertwines the
three Gauss--Manin derivatives, and the base-germ isomorphism transports
directions through \(A\).  Pairing the third derivative with the original
Hodge generator gives \cref{eq:8.2}; the two choices of generator contribute
only a common nonzero scalar.  This is the functoriality of the generalized
Yukawa tensor used for manifolds of Calabi--Yau Hodge type
\cite[\S2]{IlievManivel2015}.
\end{proof}

\begin{remark}[Necessity is not sufficiency]
Projective equality of one cubic at one point does not identify the higher
Yukawa jets, the full Gauss--Manin connection, monodromy, the integral
lattice, or an algebraic correspondence.  A central-fiber Hodge
correspondence is likewise insufficient: an algebraic correspondence would
support \cref{thm:realization-gate} only if it extended relatively and
induced the horizontal isomorphism \cref{eq:8.1}.  No Hodge, Yukawa, or
finite-prime match alone proves an isomorphism of motives.
\end{remark}

This gives a useful hierarchy of tests:
\[
\begin{array}{c}
\text{projective cubic}\\
\Downarrow\\
\text{higher jets and full connection}\\
\Downarrow\\
\text{monodromy and integral structure}\\
\Downarrow\\
\text{relative algebraic correspondence}.
\end{array}
\]
Failure at any gate rules out every stronger claim below it.  Passage
downward from a weaker match to a stronger conclusion requires new evidence.
